% Options for packages loaded elsewhere
\PassOptionsToPackage{unicode}{hyperref}
\PassOptionsToPackage{hyphens}{url}
\documentclass[
  a4paper,11pt]{article}
\usepackage{xcolor}
\usepackage[margin=1in]{geometry}
\usepackage{amsmath,amssymb}
\usepackage{graphicx}
\setcounter{secnumdepth}{-\maxdimen} % remove section numbering
\usepackage{iftex}
\ifPDFTeX
  \usepackage[T1]{fontenc}
  \usepackage[utf8]{inputenc}
  \usepackage{textcomp}
\else
  \usepackage{unicode-math}
  \defaultfontfeatures{Scale=MatchLowercase}
  \defaultfontfeatures[\rmfamily]{Ligatures=TeX,Scale=1}
\fi
\usepackage{lmodern}
\ifPDFTeX\else
  \setmainfont[]{Times New Roman}
\fi
\IfFileExists{upquote.sty}{\usepackage{upquote}}{}
\IfFileExists{microtype.sty}{%
  \usepackage[]{microtype}
  \UseMicrotypeSet[protrusion]{basicmath}
}{}
\makeatletter
\@ifundefined{KOMAClassName}{%
  \IfFileExists{parskip.sty}{%
    \usepackage{parskip}
  }{%
    \setlength{\parindent}{0pt}
    \setlength{\parskip}{6pt plus 2pt minus 1pt}}
}{% if KOMA class
  \KOMAoptions{parskip=half}}
\makeatother
\usepackage{longtable,booktabs,array}
\newcounter{none}
\usepackage{calc}
\usepackage{etoolbox}
\makeatletter
\patchcmd\longtable{\par}{\if@noskipsec\mbox{}\fi\par}{}{}
\makeatother
\IfFileExists{footnotehyper.sty}{\usepackage{footnotehyper}}{\usepackage{footnote}}
\makesavenoteenv{longtable}
\usepackage{bookmark}
\IfFileExists{xurl.sty}{\usepackage{xurl}}{}
\urlstyle{same}
\usepackage{hyperref}
\hypersetup{
  hidelinks,
  pdfcreator={LaTeX via pandoc}}
\setlength{\emergencystretch}{3em}
\providecommand{\tightlist}{%
  \setlength{\itemsep}{0pt}\setlength{\parskip}{0pt}}
\usepackage{amsthm}
\newtheorem{theorem}{Theorem}[section]
\newtheorem{proposition}[theorem]{Proposition}
\newtheorem{lemma}[theorem]{Lemma}
\newtheorem{corollary}[theorem]{Corollary}
\theoremstyle{definition}
\newtheorem{definition}[theorem]{Definition}
\theoremstyle{remark}
\newtheorem{remark}[theorem]{Remark}
\newtheorem{observation}[theorem]{Observation}
\begin{document}
\section{Formulation of Signed Area --- Derivation of Charge Structure
and Consequences of Spin-2
Configurations}\label{formulation-of-signed-area-derivation-of-charge-structure-and-consequences-of-spin-2-configurations}

\subsection{\texorpdfstring{--- Three Consequences Derived from the Sign
Product \(P\) of {[}4{]} §7
---}{--- Three Consequences Derived from the Sign Product P of {[}4{]} §7 ---}}\label{three-consequences-derived-from-the-sign-product-p-of-4-7}

\textbf{Author}: Noriaki Kihara (WF System Co., Ltd.)

\textbf{ORCID}: 0009-0006-7814-2793

\textbf{Date}: April 2026

\textbf{DOI}: https://doi.org/10.5281/zenodo.19731600

\begin{center}\rule{0.5\linewidth}{0.5pt}\end{center}

\subsection{§0 Purpose of This Paper}\label{purpose-of-this-paper}

In the remark of {[}4{]} §7, it was stated that ``the orientation
structure of the signed area \(M(\sigma)\) for square faces containing
the \(R\)-axis, and the rigorous definition of conserved quantities
based thereon, lie beyond the scope of this paper and will be addressed
in a separate work.''

This paper responds to that declaration and derives the following three
results.

\begin{enumerate}
\def\labelenumi{\arabic{enumi}.}
\tightlist
\item
  The rigorous definition of the signed area \(M\) and its relation to
  the sign product \(P\) of {[}4{]} Proposition 7.1 (§1)
\item
  The definition of the fermion area \(M_F\) and its properties
  concerning antiparticles and color charge (§2)
\item
  Consequences of the \(P = -1\) configuration in spin-2 \(tR\)-type
  (§3)
\end{enumerate}

§1--§3 are self-contained within the combinatorial framework of {[}4{]}.
§4 adopts the same stance as {[}4{]} §9 as an interpretive example,
showing the correspondence between \(M_F\) and charge quantization.

\textbf{Referenced papers}: - {[}3{]} Full Spherical Coverage of Central
Projection in Discrete Space and the Number-Theoretic Necessity of
Five-Dimensional Background Space (W5) - {[}4{]} Set Structure of the
Six-Dimensional Hyperrectangle and Its Combinatorial Properties (W8) -
{[}5{]} Interactions of Shape-Invariant Waves (W11)

\begin{center}\rule{0.5\linewidth}{0.5pt}\end{center}

\subsection{§1 Definition of Signed
Area}\label{definition-of-signed-area}

\subsubsection{1.1 2-Faces}\label{faces}

We use the set signs from {[}4{]} §1.3.3.

\textbf{Definition 1.1 (2-face)}. Among the six axes
\((x, y, z, t, R, Q)\), a set in which exactly two axes \(i, j\) have
nonzero signs (\(\sigma_i, \sigma_j \in \{+1, -1\}\)) and all remaining
axes have zero signs is called the \textbf{2-face} formed by axes
\((i, j)\).

\subsubsection{1.2 Signed Area}\label{signed-area}

By {[}4{]} §1.1, each coordinate can take nonzero values
\(v_i, v_j \neq 0\) (either real-valued or integer-valued).

\textbf{Definition 1.2 (Signed area)}. For the 2-face formed by axes
\((i, j)\), the \textbf{signed area} \(M(i, j)\) is defined as:

\[M(i, j) = v_i \cdot v_j\]

\(M\) is real-valued and carries two pieces of information: its
\textbf{sign} (orientation) and its \textbf{absolute value} (magnitude
of area):

\[\text{sign}(M) = \text{sign}(v_i) \cdot \text{sign}(v_j) \in \{+1, -1\}\]

\[|M| = |v_i| \cdot |v_j|\]

\(\text{sign}(M) = +1\) is called \textbf{positive orientation} and
\(\text{sign}(M) = -1\) is called \textbf{reverse orientation}. The
specific values of \(|v_i|\) and \(|v_j|\) are not constrained by this
definition.

\textbf{Remark}. For cancellation of signed areas (e.g.,
particle-antiparticle pair annihilation) to satisfy \(M + M' = 0\), the
condition \(|M| = |M'|\), i.e., equality of the absolute values of the
axes, is required. The antiparticle definition in {[}5{]} §5.4 (only
spatial axis signs are inverted; absolute values are invariant)
automatically satisfies this condition.

\textbf{Proposition 1.2 (Signed area on integer lattices)}. By {[}3{]},
coordinates that realize full spherical coverage of central projection
in discrete space take values on the integer lattice \(\mathbb{Z}^n\).
Then for nonzero coordinates
\(v_i, v_j \in \mathbb{Z} \setminus \{0\}\):

\begin{enumerate}
\def\labelenumi{(\roman{enumi})}
\item
  \(M(i,j) = v_i \cdot v_j \in \mathbb{Z} \setminus \{0\}\)
\item
  \(\text{sign}(M) \in \{+1, -1\}\)
\end{enumerate}

\emph{Proof}. (i) The integer ring \(\mathbb{Z}\) is closed under
multiplication (\(a, b \in \mathbb{Z} \implies ab \in \mathbb{Z}\)) and
is an integral domain (\(a \neq 0\) and
\(b \neq 0 \implies ab \neq 0\)). Since
\(v_i, v_j \in \mathbb{Z} \setminus \{0\}\), we have
\(M = v_i \cdot v_j \in \mathbb{Z} \setminus \{0\}\). (ii) Since
\(M \neq 0\), \(\text{sign}(M) \in \{+1, -1\}\) is well-defined.
\(\square\)

\textbf{Remark}. Definition 1.2 is also applicable to the continuous
case where coordinate values are real-valued, but the two-valuedness of
\(\text{sign}(M)\) is guaranteed solely by the condition
\(v_i, v_j \neq 0\). On integer lattices, this condition is
automatically satisfied by the lattice structure.

\subsubsection{\texorpdfstring{1.3 Relation to the Sign Product
\(P\)}{1.3 Relation to the Sign Product P}}\label{relation-to-the-sign-product-p}

\textbf{Proposition 1.1}. The sign product
\(P = (-1)^{|\{i \in \{t, R\} : \sigma_i < 0\}|}\) from {[}4{]}
Proposition 7.1 coincides with the sign of \(M(t, R)\) when both the
\(t\)-axis and \(R\)-axis are nonzero: \(P = \text{sign}(M(t, R))\).

\emph{Proof}. When
\(\text{sign}(v_t), \text{sign}(v_R) \in \{+1, -1\}\):

{\def\LTcaptype{none} % do not increment counter
\begin{longtable}[]{@{}
  >{\centering\arraybackslash}p{(\linewidth - 8\tabcolsep) * \real{0.2000}}
  >{\centering\arraybackslash}p{(\linewidth - 8\tabcolsep) * \real{0.2000}}
  >{\centering\arraybackslash}p{(\linewidth - 8\tabcolsep) * \real{0.2000}}
  >{\centering\arraybackslash}p{(\linewidth - 8\tabcolsep) * \real{0.2000}}
  >{\centering\arraybackslash}p{(\linewidth - 8\tabcolsep) * \real{0.2000}}@{}}
\toprule\noalign{}
\begin{minipage}[b]{\linewidth}\centering
\(\text{sign}(v_t)\)
\end{minipage} & \begin{minipage}[b]{\linewidth}\centering
\(\text{sign}(v_R)\)
\end{minipage} & \begin{minipage}[b]{\linewidth}\centering
Number of negatives
\end{minipage} & \begin{minipage}[b]{\linewidth}\centering
\(P\)
\end{minipage} & \begin{minipage}[b]{\linewidth}\centering
\(\text{sign}(M)\)
\end{minipage} \\
\midrule\noalign{}
\endhead
\bottomrule\noalign{}
\endlastfoot
\(+\) & \(+\) & 0 & \(+1\) & \(+1\) \\
\(+\) & \(-\) & 1 & \(-1\) & \(-1\) \\
\(-\) & \(+\) & 1 & \(-1\) & \(-1\) \\
\(-\) & \(-\) & 2 & \(+1\) & \(+1\) \\
\end{longtable}
}

\(P = \text{sign}(M)\) holds for all four configurations. \(\square\)

\textbf{Remark}. When only one of \(t, R\) is nonzero,
\(P = (-1)^{0 \text{ or } 1}\), and the 2-face \(M(t,R)\) is not defined
(a single axis does not constitute a face). In this case, \(P\)
functions as a generalization of \(\text{sign}(M)\).

\begin{center}\rule{0.5\linewidth}{0.5pt}\end{center}

\subsection{§2 Fermion Area}\label{fermion-area}

\subsubsection{2.1 Definition}\label{definition}

By {[}4{]} Definition 2.1, fermion types have exactly \(k = 2\) of the
four \(xyzt\) axes nonzero.

\textbf{Definition 2.1 (Fermion area)}. The signed area
\(M_F = v_i \cdot v_j\) of the 2-face formed by the two nonzero axes
\((i, j) \subseteq \{x, y, z, t\}\) of a fermion is called the
\textbf{fermion area}. Its \textbf{orientation} is given by
\(\text{sign}(M_F) = \text{sign}(v_i) \cdot \text{sign}(v_j) \in \{+1, -1\}\).

\subsubsection{2.2 Orientation of Fermion Area in {[}4{]} Table
A}\label{orientation-of-fermion-area-in-4-table-a}

We compute \(\text{sign}(M_F)\) for all fermions in {[}4{]} §9.9 Table
A. The absolute values \(|v_i|\) of each axis do not affect the results
of this section.

\textbf{Charged leptons} (nonzero axes: 1 spatial axis + \(t\)-axis,
\(k_t < 0\) (down-type), \(Q = 0\)):

{\def\LTcaptype{none} % do not increment counter
\begin{longtable}[]{@{}llcc@{}}
\toprule\noalign{}
Particle & Set & Nonzero axes & \(\text{sign}(M_F)\) \\
\midrule\noalign{}
\endhead
\bottomrule\noalign{}
\endlastfoot
\(e^-\) & \((+, 0, 0, -, 0, 0)\) & \(x, t\) & \((+)(-)= -1\) \\
\(\mu^-\) & \((0, +, 0, -, 0, 0)\) & \(y, t\) & \((+)(-)= -1\) \\
\(\tau^-\) & \((0, 0, +, -, 0, 0)\) & \(z, t\) & \((+)(-)= -1\) \\
\end{longtable}
}

\textbf{Neutrinos} (nonzero axes: 2 spatial axes, \(t = 0\)):

{\def\LTcaptype{none} % do not increment counter
\begin{longtable}[]{@{}llcc@{}}
\toprule\noalign{}
Particle & Set & Nonzero axes & \(\text{sign}(M_F)\) \\
\midrule\noalign{}
\endhead
\bottomrule\noalign{}
\endlastfoot
\(\nu_e\) & \((+, +, 0, 0, 0, 0)\) & \(x, y\) & \((+)(+)= +1\) \\
\(\nu_\mu\) & \((0, +, +, 0, 0, 0)\) & \(y, z\) & \((+)(+)= +1\) \\
\(\nu_\tau\) & \((+, 0, +, 0, 0, 0)\) & \(x, z\) & \((+)(+)= +1\) \\
\end{longtable}
}

\textbf{Up-type quarks} (nonzero axes: 1 spatial axis + \(t\)-axis,
\(k_t > 0\) (up-type), \(Q \in \{1, 2, 4\}\)):

{\def\LTcaptype{none} % do not increment counter
\begin{longtable}[]{@{}llcc@{}}
\toprule\noalign{}
Particle & Set & Nonzero axes & \(\text{sign}(M_F)\) \\
\midrule\noalign{}
\endhead
\bottomrule\noalign{}
\endlastfoot
\(u\) & \((+, 0, 0, +, 0, Q)\) & \(x, t\) & \((+)(+)= +1\) \\
\(c\) & \((0, +, 0, +, 0, Q)\) & \(y, t\) & \((+)(+)= +1\) \\
\(t\) & \((0, 0, +, +, 0, Q)\) & \(z, t\) & \((+)(+)= +1\) \\
\end{longtable}
}

\textbf{Down-type quarks} (nonzero axes: 1 spatial axis + \(t\)-axis,
\(k_t < 0\) (down-type), \(Q \in \{1, 2, 4\}\)):

{\def\LTcaptype{none} % do not increment counter
\begin{longtable}[]{@{}llcc@{}}
\toprule\noalign{}
Particle & Set & Nonzero axes & \(\text{sign}(M_F)\) \\
\midrule\noalign{}
\endhead
\bottomrule\noalign{}
\endlastfoot
\(d\) & \((+, 0, 0, -, 0, Q)\) & \(x, t\) & \((+)(-)= -1\) \\
\(s\) & \((0, +, 0, -, 0, Q)\) & \(y, t\) & \((+)(-)= -1\) \\
\(b\) & \((0, 0, +, -, 0, Q)\) & \(z, t\) & \((+)(-)= -1\) \\
\end{longtable}
}

\textbf{Summary}: It is solely the \textbf{orientation} (sign) of
\(M_F\) that determines the charge structure; it does not depend on the
absolute value \(|M_F| = |v_i| \cdot |v_j|\).

{\def\LTcaptype{none} % do not increment counter
\begin{longtable}[]{@{}llcc@{}}
\toprule\noalign{}
Fermion type & Type of nonzero axes & \(Q\) & \(\text{sign}(M_F)\) \\
\midrule\noalign{}
\endhead
\bottomrule\noalign{}
\endlastfoot
Neutrino & (spatial, spatial) & 0 & \(+1\) \\
Charged lepton & (spatial, \(t<0\)) & 0 & \(-1\) \\
Up-type quark & (spatial, \(t>0\)) & 1,2,4 & \(+1\) \\
Down-type quark & (spatial, \(t<0\)) & 1,2,4 & \(-1\) \\
\end{longtable}
}

\subsubsection{2.3 Fermion Area of
Antiparticles}\label{fermion-area-of-antiparticles}

By {[}4{]} §9.1, for antiparticles the signs of both the spatial axes
\((x, y, z)\) and the \(t\)-axis are inverted, and all \(Q\) bits are
inverted (\((R,G,B) \to (\bar{R},\bar{G},\bar{B})\)).

\textbf{Proposition 2.1}. When the fermion face consists of a (spatial,
\(t\)) combination, the fermion area of the antiparticle is identical to
that of the particle: \(M_F^{\text{anti}} = M_F\).

\emph{Proof}. For the antiparticle, the spatial axis value is inverted
to \(-v_{\text{spatial}}\) and the \(t\)-axis value is inverted to
\(-v_t\). Therefore,
\(M_F^{\text{anti}} = (-v_{\text{spatial}}) \cdot (-v_t) = v_{\text{spatial}} \cdot v_t = M_F\).
\(\square\)

\textbf{Proposition 2.2}. When the fermion face consists of a (spatial,
spatial) combination, the fermion area of the antiparticle is identical
to that of the particle: \(M_F^{\text{anti}} = M_F\).

\emph{Proof}. For the antiparticle, both spatial axis values are
inverted to \(-v_{\text{spatial}_1}, -v_{\text{spatial}_2}\). Therefore,
\(M_F^{\text{anti}} = (-v_{\text{spatial}_1}) \cdot (-v_{\text{spatial}_2}) = v_{\text{spatial}_1} \cdot v_{\text{spatial}_2} = M_F\).
\(\square\)

\textbf{Corollary 2.1}. For all fermion types,
\(M_F^{\text{anti}} = M_F\) holds. The fermion area is invariant under
the antiparticle transformation.

\subsubsection{2.4 Independence of Color Charge and Fermion
Area}\label{independence-of-color-charge-and-fermion-area}

\textbf{Proposition 2.3}. The value of the \(Q\)-axis does not affect
the fermion area \(M_F\). Fermions with the same
\(\sigma_{\text{spatial}}, \sigma_t\) have identical \(M_F\) regardless
of the value of \(Q\).

\emph{Proof}.
\(\text{sign}(M_F) = \text{sign}(v_i) \cdot \text{sign}(v_j)\) is
determined solely by the values of the \(xyzt\) axes, and \(Q\) does not
appear in the definition. \(\square\)

\begin{center}\rule{0.5\linewidth}{0.5pt}\end{center}

\subsection{\texorpdfstring{§3 The \(P = -1\) Configuration of Spin-2
\(tR\)-Type}{§3 The P = -1 Configuration of Spin-2 tR-Type}}\label{the-p--1-configuration-of-spin-2-tr-type}

By {[}4{]} Proposition 7.1, spin-2 \(tR\)-type admits both \(P = +1\)
and \(P = -1\) configurations. In {[}4{]} §9.7, the physical
correspondence of \(P = -1\) was left unresolved.

\textbf{Proposition 3.1}. The \(P = -1\) configuration of spin-2
\(tR\)-type is equivalent to \(\sigma_t\) and \(\sigma_R\) having
opposite signs.

\emph{Proof}. By Proposition 1.1,
\(P = \text{sign}(M(t, R)) = \text{sign}(v_t) \cdot \text{sign}(v_R)\).
\(P = -1\) \(\iff\) \(\text{sign}(v_t) \cdot \text{sign}(v_R) = -1\)
\(\iff\) \(v_t\) and \(v_R\) have opposite signs. \(\square\)

\textbf{Proposition 3.2}. Applying the sign rule for
attraction/repulsion from {[}5{]} §4.3 (same relative sign of \(t\)-axis
components implies attraction; opposite signs imply repulsion) to the
graviton interpretation of {[}4{]} §9.5:

\begin{itemize}
\tightlist
\item
  \(P = +1\) (\(t, R\) same sign): \(t\)-component has a single sign
  \(\to\) attraction only
\item
  \(P = -1\) (\(t, R\) opposite signs): \(t\)-component and
  \(R\)-component have opposite signs \(\to\) repulsive along the
  \(R\)-axis direction
\end{itemize}

\emph{Proof}. In {[}5{]} §4.3, the sign of interaction between
direction-constrained waves is determined by
\(\text{sign}(k_t^{(1)}) \cdot \text{sign}(k_t^{(2)})\). When the
spherical wave of spin-2 \(tR\)-type follows this sign rule, \(P = +1\)
implies both axes share the same directionality, yielding attraction
only. For \(P = -1\), the action along the \(R\)-direction has the
opposite sign to the \(t\)-direction, acting repulsively with respect to
the spatial scale. \(\square\)

\textbf{Remark}. The interpretation of the physical consequences of the
\(P = -1\) configuration is addressed in §4.3.

\begin{center}\rule{0.5\linewidth}{0.5pt}\end{center}

\subsection{§4 Interpretive Example: Derivation of Charge
Structure}\label{interpretive-example-derivation-of-charge-structure}

This section has the same status as {[}4{]} §9; it is an interpretive
example that maps the combinatorial results of §1--§3 onto the Standard
Model.

\subsubsection{\texorpdfstring{4.1 Third Component of Weak Isospin
\(I_3\)}{4.1 Third Component of Weak Isospin I\_3}}\label{third-component-of-weak-isospin-i_3}

Using the \(\text{sign}(M_F)\) introduced in §2.2, we define the third
component of weak isospin \(I_3\) of the Standard Model as follows.

In {[}4{]} §9.1, weak isospin is encoded by the sign of the \(t\)-axis
(\(k_t > 0\): up-type, \(k_t < 0\): down-type). Since charged leptons
are down-type (\(k_t < 0\)), they have \(\text{sign}(M_F) = -1\), and
\(\text{sign}(M_F)\) directly distinguishes up-type (\(+1\)) from
down-type (\(-1\)) for all fermions.

\[I_3 = \frac{\text{sign}(M_F)}{2}\]

{\def\LTcaptype{none} % do not increment counter
\begin{longtable}[]{@{}lcccc@{}}
\toprule\noalign{}
Fermion type & \(Q\) & \(\text{sign}(M_F)\) & \(I_3\) & SM \(I_3\) \\
\midrule\noalign{}
\endhead
\bottomrule\noalign{}
\endlastfoot
Neutrino & \(0\) & \(+1\) & \(+1/2\) & \(+1/2\) \\
Charged lepton & \(0\) & \(-1\) & \(-1/2\) & \(-1/2\) \\
Up-type quark & \(1,2,4\) & \(+1\) & \(+1/2\) & \(+1/2\) \\
Down-type quark & \(1,2,4\) & \(-1\) & \(-1/2\) & \(-1/2\) \\
\end{longtable}
}

The value of \(I_3 = \text{sign}(M_F)/2\) agrees with the Standard Model
for all fermions. Since isospin is encoded by the \(t\)-axis sign, a
single unified formula applies to both quarks and leptons.

\subsubsection{\texorpdfstring{4.2 Determination of Weak Hypercharge
\(Y\) from the \(Q\)
Value}{4.2 Determination of Weak Hypercharge Y from the Q Value}}\label{determination-of-weak-hypercharge-y-from-the-q-value}

The Standard Model charge formula (Gell-Mann--Nishijima relation):

\[Q_{\text{charge}} = I_3 + \frac{Y}{2}\]

Assume that \(Y\) depends only on the lepton/quark classification:
\(Y = Y_L\) when \(Q = 0\) (leptons, set-bit count 0), and \(Y = Y_Q\)
when \(Q \in \{1,2,4\}\) (quarks, set-bit count 1). Then the following
two conditions uniquely determine \(Y_L, Y_Q\).

\textbf{Condition (i): Integer charges for leptons}. For any fermion
with \(Q = 0\), \(Q_{\text{charge}} = I_3 + Y_L/2 \in \mathbb{Z}\) must
hold. Since \(I_3 \in \{+1/2, -1/2\}\), \(Y_L\) must be odd.

\textbf{Condition (ii): Integer charges for color-neutral baryons}. The
total charge of a three-body color-neutral system (three single-bit
colors \(Q = 1, 2, 4\), one each),
\(Q_{\text{baryon}} = \sum_{a=1}^{3}(I_3^{(a)} + Y_Q / 2) \in \mathbb{Z}\),
must hold for any combination of \(I_3^{(a)} \in \{+1/2, -1/2\}\).

\[Q_{\text{baryon}} = \sum I_3^{(a)} + \frac{3 Y_Q}{2}\]

Since \(\sum I_3^{(a)} \in \{-3/2, -1/2, +1/2, +3/2\}\) (always a
half-integer), \(3 Y_Q / 2\) must also be a half-integer. That is,
\(Y_Q = (2n+1)/3\) (\(n \in \mathbb{Z}\)).

\textbf{Determination of the minimal solution}. From condition (i), the
odd number with smallest \(|Y_L|\) is \(Y_L = \pm 1\). Requiring
negative charge for charged leptons:

\[Q_{\text{charge}}(e^-) = -\frac{1}{2} + \frac{Y_L}{2} < 0 \implies Y_L < 1\]

Hence \(Y_L = -1\). From condition (ii), the solution with smallest
\(|Y_Q|\) is \(n = 0\): \(Y_Q = 1/3\).

\textbf{Proposition 4.1 (Unique determination of \(Y\) from integer
charge conditions)}. From the two-valuedness of
\(I_3 \in \{+1/2, -1/2\}\) and the condition that charges of observable
particles (leptons and color-neutral baryons) are integers, the minimal
solution for \(Y\) is uniquely determined as:

\[Y_L = -1 \quad (Q = 0), \qquad Y_Q = +\frac{1}{3} \quad (Q \in \{1,2,4\})\]

\subsubsection{4.3 Consequences}\label{consequences}

\textbf{Consequence 1: Reproduction of the Gell-Mann--Nishijima
formula}. Substituting the \(I_3\) from §4.1 and \(Y\) from Proposition
4.1:

{\def\LTcaptype{none} % do not increment counter
\begin{longtable}[]{@{}
  >{\raggedright\arraybackslash}p{(\linewidth - 10\tabcolsep) * \real{0.1379}}
  >{\centering\arraybackslash}p{(\linewidth - 10\tabcolsep) * \real{0.1724}}
  >{\centering\arraybackslash}p{(\linewidth - 10\tabcolsep) * \real{0.1724}}
  >{\centering\arraybackslash}p{(\linewidth - 10\tabcolsep) * \real{0.1724}}
  >{\centering\arraybackslash}p{(\linewidth - 10\tabcolsep) * \real{0.1724}}
  >{\centering\arraybackslash}p{(\linewidth - 10\tabcolsep) * \real{0.1724}}@{}}
\toprule\noalign{}
\begin{minipage}[b]{\linewidth}\raggedright
Fermion
\end{minipage} & \begin{minipage}[b]{\linewidth}\centering
\(\text{sign}(M_F)\)
\end{minipage} & \begin{minipage}[b]{\linewidth}\centering
\(I_3\)
\end{minipage} & \begin{minipage}[b]{\linewidth}\centering
\(Y\)
\end{minipage} & \begin{minipage}[b]{\linewidth}\centering
\(Q_{\text{charge}} = I_3 + Y/2\)
\end{minipage} & \begin{minipage}[b]{\linewidth}\centering
SM value
\end{minipage} \\
\midrule\noalign{}
\endhead
\bottomrule\noalign{}
\endlastfoot
\(\nu\) & \(+1\) & \(+1/2\) & \(-1\) & \(0\) & \(0\) ✓ \\
\(e^-\) & \(-1\) & \(-1/2\) & \(-1\) & \(-1\) & \(-1\) ✓ \\
\(u\) & \(+1\) & \(+1/2\) & \(+1/3\) & \(+2/3\) & \(+2/3\) ✓ \\
\(d\) & \(-1\) & \(-1/2\) & \(+1/3\) & \(-1/3\) & \(-1/3\) ✓ \\
\end{longtable}
}

This is consistent with the charge map \(f_Q^{\text{charge}}\) of
{[}4{]} Definition 9.1.

\textbf{Consequence 2: Automatic satisfaction of anomaly cancellation}.
The solution from Proposition 4.1 satisfies:

\[Y_L + 3 Y_Q = -1 + 3 \times \frac{1}{3} = 0\]

This corresponds to the gauge anomaly cancellation condition in the
Standard Model. In this framework, it is derived algebraically from the
set-bit count structure of \(Q\) (one color-neutral state \(Q = 0\) +
three color charge states \(Q \in \{1,2,4\}\)) and the integer charge
condition.

\textbf{Consequence 3: Charge equality of proton and electron}. The
charge of the proton \((uud)\):

\[Q_p = \sum I_3^{(a)} + \frac{3 Y_Q}{2} = \left(\frac{1}{2} + \frac{1}{2} - \frac{1}{2}\right) + \frac{1}{2} = +1\]

The charge (absolute value) of the electron:

\[|Q_e| = \left|I_3(e^-) + \frac{Y_L}{2}\right| = \left|-\frac{1}{2} - \frac{1}{2}\right| = 1\]

\(|Q_p| = |Q_e| = 1\); the reason that the proton and electron carry
charges of equal magnitude is determined algebraically from the set-bit
count structure of \(Q\) (\(1 + 3\) classification: 1 lepton state + 3
quark colors) and the two-valuedness of \(I_3\) (\(\pm 1/2\)). There is
no need to introduce fractional charges \(+2/3, -1/3\) as fundamental
quantities.

\textbf{Consequence 4: Interpretation of spin-2 \(P = -1\)}. By
Proposition 3.2 of §3, the graviton configuration with \(P = -1\) acts
repulsively along the \(R\)-axis direction. Since the \(R\)-axis in
{[}4{]} §9.8 is the axis that links the subjective coordinate system to
the spatial scale of the six-dimensional structure, the \(P = -1\)
configuration may correspond to \textbf{accelerating expansion of the
spatial scale}. This is a candidate for the geometric description of the
cosmological constant \(\Lambda > 0\) (dark energy). However, whether
the \(P = -1\) configuration is physically realized (existence of
selection rules) remains an open question, and this paper limits itself
to stating the consequence.

\begin{center}\rule{0.5\linewidth}{0.5pt}\end{center}

\subsection{§5 Conclusion}\label{conclusion}

In this paper, we provided the definition of the signed area
\(M(\sigma)\) declared in {[}4{]} §7 and derived the following three
results.

\begin{enumerate}
\def\labelenumi{\arabic{enumi}.}
\tightlist
\item
  The \textbf{signed area \(M(i,j) = v_i \cdot v_j\)} (Definition 1.2)
  coincides with the restriction of the sign product \(P\) of {[}4{]}
  Proposition 7.1 to the case where both \(t\) and \(R\) are nonzero
  (Proposition 1.1).
\item
  The \textbf{fermion area \(M_F\)} (Definition 2.1) is computable for
  all fermions in {[}4{]} Table A, and we demonstrated invariance under
  the antiparticle transformation for all fermion types (Propositions
  2.1, 2.2), and independence from the \(Q\)-axis (Proposition 2.3).
\item
  The \textbf{spin-2 \(P = -1\)} configuration corresponds to opposite
  signs of \(v_t, v_R\), and from the sign rule of {[}5{]}, it was
  derived to be repulsive along the \(R\)-axis direction (Propositions
  3.1, 3.2).
\end{enumerate}

In the interpretive example (§4): - \(I_3 = \text{sign}(M_F)/2\) agrees
uniformly with the third component of weak isospin in the Standard Model
for all fermions (§4.1) - The integer charge condition uniquely
determines \(Y_L = -1, Y_Q = +1/3\) (Proposition 4.1) - The anomaly
cancellation condition \(Y_L + 3Y_Q = 0\) is automatically satisfied
(Consequence 2) - The charge equality of the proton and electron is
derived from the \(1+3\) set-bit count structure of \(Q\) (Consequence
3)

were demonstrated.

\begin{center}\rule{0.5\linewidth}{0.5pt}\end{center}

\subsection{References}\label{references}

\begin{itemize}
\tightlist
\item
  {[}3{]} N. Kihara, ``Full Spherical Coverage of Central Projection in
  Discrete Space and the Number-Theoretic Necessity of Five-Dimensional
  Background Space,'' Zenodo preprint, DOI: 10.5281/zenodo.19624957
  (2026).
\item
  {[}4{]} N. Kihara, ``Set Structure of the Six-Dimensional
  Hyperrectangle and Its Combinatorial Properties,'' Zenodo preprint,
  DOI: 10.5281/zenodo.19721125 (2026).
\item
  {[}5{]} N. Kihara, ``Interactions of Shape-Invariant Waves --- Axial
  Displacement Transfer, Wave Packet Deformation, and Retroactive
  Construction of Causality,'' Zenodo preprint, DOI:
  10.5281/zenodo.19721128 (2026).
\end{itemize}
\end{document}
